% ------------------------------------------------------------------------
% file `3_uaa1_fonctions_2023-2024-connaitre-4-solution-body.tex'
%   in folder `xsim/'
%
%     solution of type `connaitre' with id `4'
%
% generated by the `solconnaitre' environment of the
%   `xsim' package v0.21 (2022/02/12)
% from source `3_uaa1_fonctions_2023-2024' on 2023/10/18 on line 67
% ------------------------------------------------------------------------
    \begin{tasks}
        \task Faux. Les deux membres ont été divisés par \(-3\), il faut donc changer le sens de l'inégalité.
        \task Vrai. Si \(x=0\) alors \(f(0) = 2\cdot 0 + 8 = 8\).
        \task Faux. Les deux droites sont confondues.
        \task Faux. Cela dépend du signe de son ordonnée à l'origine, si elle est positive alors la racine est négative sinon la racine est positive.
    \end{tasks}
